\documentclass[10pt,a4paper]{amsart}
\usepackage[margin=19mm]{geometry}
\usepackage{amsmath,amssymb,mathtools}
\usepackage[scr=rsfs,cal=euler]{mathalfa}
\usepackage{xfrac}
\usepackage[unicode,hidelinks,bookmarks=false]{hyperref}
\newcommand{\ie}{i.e.\ }
\newcommand{\cf}{cf.\ }
\newcommand{\resp}{respectively\ }
\newcommand{\Subsection}{\subsection}
\providecommand{\nobreakdash}{\nobreak\hbox{-}}
\setlength{\emergencystretch}{2em}
\multlinegap0pt
\usepackage{amssymb}
\usepackage{algorithm}\usepackage{algpseudocode}
\usepackage[shortlabels]{enumitem}
\newtheorem{thm}{Theorem}
\newcommand{\Zn}[1]{\mathbb{Z}_{#1}}
\newcommand{\supp}{\mathrm{supp}}
\newcommand{\wt}{\mathrm{wt}}
\title{Construction of Multicyclic Codes of Arbitrary Dimension $r$ via Idempotents: A Unified Combinatorial-Algebraic Approach}
\author{Jean Charles Ramanandraibe, Ramamonjy Andriamifidisoa}
\date{Source: arXiv:2603.07177v1 (CC BY 4.0)}
\begin{document}
\maketitle

\noindent\emph{Source and licence.} Construction of Multicyclic Codes of Arbitrary Dimension $r$ via Idempotents: A Unified Combinatorial-Algebraic Approach. Version arXiv:2603.07177v1 (CC BY 4.0), distributed by arXiv under the Creative Commons Attribution 4.0 International licence, http://creativecommons.org/licenses/by/4.0/, as recorded in the arXiv licence metadata for this exact version and shown on the abstract page https://arxiv.org/abs/2603.07177v1. Full licence notice: \texttt{licenses/arxiv-2603.07177-CC-BY-4.0.txt}.\par\smallskip

\noindent\textit{Editorial scope.} Counted unit: the article's thm thm:bound (source0.tex lines 211--216). The article's complete original proof of it is delivered with it (source0.tex lines 217--253, one \texttt{proof} environment of 37 lines). No mathematics is written, repaired or re-derived: the statement and the proof are the article's own lines, in the article's own notation, and no retained line is substituted or re-flowed.  Retained uncounted as the source-local context the proof needs: lines 183--189.  Nothing was omitted from any retained span, so the package carries no omission record.  No retained line contains a citation, so the wrapper carries no bibliography and no bibliography entry.  No retained line contains a figure environment, an image-inclusion command or drawing code, so no image is delivered and the package declares no asset dependency.  Editorial formatting only, in addition: an AMS \texttt{amsart} preamble that restates the article's own declarations and macros used by the retained lines, namely the environments \texttt{thm} on separate counters, together with the standard packages those lines need.  The retained environments print in this document as Theorem 1 in order of appearance, whereas the article prints the same items with the numbering of its own full document.  The complete native source of the article as distributed in the arXiv e-print for this exact version is bundled unchanged as \texttt{sources/195922/source0.tex}.
\begin{thm}[Polynomial basis]\label{thm:basis}
Let $k_t$ be minimal such that $X_t^{k_t} e \in \mathrm{span}\{X_t^m e : m < k_t\}$. Then
\[
\mathcal{B} = \{ X_1^{m_1}\cdots X_r^{m_r} e : 0\le m_t < k_t \}
\]
is a basis and $\dim C = \prod k_t$.
\end{thm}

\begin{thm}[Product bound]\label{thm:bound}
For $C = \langle e \rangle$,
\[
n = \prod n_t, \quad k = \prod k_t, \quad d \ge \prod_{t=1}^r (n_t - k_t + 1).
\]
\end{thm}

\begin{proof}
The length $n$ is the total number of monomials in $R$, so
\[
n = \prod_{t=1}^r n_t.
\]
The dimension $k$ is given by Theorem \ref{thm:basis}:
\[
k = \prod_{t=1}^r k_t.
\]
Let $c = P e \in C \setminus \{0\}$, with $\deg_{X_t} P < k_t$ for all $t$.
For each $t = 1,\dots,r$, let
\[
A_t = \{ i_t \in \Zn{n_t} \mid \exists (i_1,\dots,i_r) \in \supp(c)
\text{ with this coordinate } i_t \}.
\]
Then
\[
\supp(c) \subseteq A_1 \times \cdots \times A_r
\quad\text{and thus}\quad
\wt(c) \ge |A_1|\cdots |A_r|.
\]
Fix $t \in \{1,\dots,r\}$ and freeze the variables $\{X_j\}_{j \neq t}$.
The word $c$ can then be seen as a non-zero word of a univariate cyclic code in the variable $X_t$, of length $n_t$ and dimension $k_t$.
By the Singleton bound (or BCH bound for cyclic codes), every non-zero word satisfies:
\[
|A_t| \ge n_t - k_t + 1.
\]
Since this estimate holds for all $t=1,\dots,r$, we obtain:
\[
\wt(c) \ge \prod_{t=1}^r |A_t|
\ge \prod_{t=1}^r (n_t - k_t + 1).
\]
Thus, the minimum distance $d$ of the code $C$ satisfies:
\[
d \ge \prod_{t=1}^r (n_t - k_t + 1).
\]
\end{proof}

\end{document}
